\section{Uvod}

Znameniti Gelfandov izrek pravi, da je vsaka komutativna $C^*$-algebra 
$*$-izomorfna $C^*$-algebri zveznih funkcij na kompaktnem Hausdorffovem prostoru $X$.
To vodi do fundamentalne ideje \emph{nekomutativne geometrije}: na splošno (ne nujno komutativno) $C^*$-algebro
lahko gledamo kot na algebro funkcij na nekomutativni posplošitvi pojma kompaktnega Hausdorffovega topološkega prostora.
Nekatere klasične geometrijske objekte, ki jih lahko formuliramo v jeziku komutativnih $C^*$-algeber $C(X)$, 
tako vpeljemo tudi v nekomutativnem kontekstu. 
Glavni cilj tega magistrskega dela je izpostaviti dva takšna geometrijska pojma, \emph{rob Šilova} in \emph{Choqueta},
in vpeljati njuna nekomutativna analoga.

Ključen pojem, ki ga bomo pri tem potrebovali, so \emph{povsem pozitivne preslikave}.
Eden izmed načinov, na katere lahko motiviramo povsem pozitivne preslikave motiviramo, pride iz nekomutativne verjetnosti.
Povsem pozitivne preslikave so namreč nekomutativna posplošitev pojma {Markovskega jedra} iz klasične verjetnosti.
Običajnemu verjetnostnemu prostoru $(X, \mathcal{F}, \mu)$ lahko priredimo $C^*$-algebro $\mathcal{B} (X)$,
ki jo sestavljajo vse omejene kompleksne naključne spremenljivke na $X$, in pričakovano vrednost $\E: \mathcal{B} (X) \to \C$.
To nam omogoča, da na verjetnostne prostore gledamo iz bolj abstraktnega vidika
kot par $(\mathfrak{A}, \E)$,
kjer je $\mathfrak{A}$ $C^*$-algebra in $\E: \mathfrak{A} \to \C$ pozitiven funkcional,
za katerega velja $\E (1) = 1$ (takšnim funkcionalom v žargonu pravimo \emph{stanja}).
Izkaže se, da lahko večino običajnih pojmov iz verjetnostne teorije, kot na primer \emph{dogodke}, \emph{slučajne spremenljivke}, ali pa \emph{pogojno pričakovano vrednost}, 
prevedemo tudi na nekomutativni kontekst
(za več o tej temi glej \cite{nicaspeicher2006}). Oglejmo si Markovsko jedro za merljiva prostora $(X, \mathcal{F})$ in $(Y, \mathcal{G})$,
ki je definirano kot funkcija $K: Y \times \mathcal{F} \to [0, 1]$,
tako da je za vsak $A \in \mathcal{F}$ preslikava $y \mapsto K(y, A)$ $\mathcal{G}$-merljiva in 
za vsak $y \in Y$ preslikava $B \mapsto K(y, B)$ verjetnostna mera na $(X, \mathcal{F})$.
To inducira preslikavo 
$$\varphi: \mathcal{B}(X) \to \mathcal{B}(Y),\quad (\varphi f)(y) = \int_X f(x)\, K(y, dx),$$
ki ima lastnost, da za poljubne funkcije $a_1, \dots, a_n \in \mathcal{B}(X)$ in $b_1, \dots, b_n \in \mathcal{B}(Y)$ velja 
\begin{equation}\label{eq:intro1}
    \sum_{i, j = 1} ^n b_i ^* \varphi(a_i ^* a_j) b_j \geq 0.
\end{equation}
Linearnim preslikavam $\varphi: \mathfrak{A} \to \mathfrak{B}$ med $C^*$-algebrami,
za katere velja lastnost \eqref{eq:intro1}, pravimo povsem pozitivne preslikave.
Povsem pozitivne preslikave imajo med drugim aplikacije v opisu \emph{odprtih kvantnih sistemov} s pomočjo
\emph{kvantnih dinamičnih polgrup} %, ki so nekomutativna posplošitev klasičnih markovskih polgrup 
(\cite{attaljoyepillet2006}),
ali pa v \emph{kvantni teoriji informacij}, kjer so znani kot \emph{kvantni kanali} (\cite{nielsenchuang2010}).

Pri organizaciji snovi v tem magistrskem delu sem želel okvirno slediti zgodovinskemu zaporedju teorije povsem pozitivnih preslikav.
Za njen začetek lahko vzamemo Stinespringov članek \cite{stinespring1955} iz leta 1955
in Arvesonov članek \cite{arveson1969} iz leta 1969 -- na slednjega se bomo sklicevali skozi celotno magistrsko delo.
Posebnega pomena sta Stinespringov izrek, ki je posplošitev Gelfand-Naimark-Segalovega izreka,
in Arvesonov razširitveni izrek, ki je posplošitev Hahn-Banachovega izreka za povsem pozitivne preslikave.
Posledica teh dveh izrekov je, da lahko povsem pozitivne preslikave $\varphi: \mathfrak{A} \to \mathfrak{B}$ smatramo za naravno posplošitev pozitivnih funkcionalov oz.~stanj
na $C^*$-algebrah, le da je kodomena $\mathfrak{B}$ sedaj poljubna $C^*$-algebra. Osnovne primere in lastnosti povsem pozitivnih preslikav 
obravnavamo v poglavju \ref{sec:2}.

Iz Arvesonovega dela je sledilo, da naravna domena za povsem pozitivne preslikave niso celotne $C^*$-algebre,
temveč njihovi enotski sebiadjungirani podprostori, ki jim v žargonu pravimo \emph{operatorski sistemi}.
Kot bomo videli, so povsem pozitivne preslikave morfizmi v kategoriji operatorskih sistemov.
Operatorski sistemi se sicer pojavijo tudi v nekomutativni geometriji%, ko obravnavamo \emph{spektralne trojice} -- to so trojice $(\mathfrak{A}, \mathcal{H}, D)$,
%kjer je $\mathfrak{A}$ enotska $*$-algebra omejenih operatorjev na Hilbertovem prostoru $\mathcal{H}$ in $D$ neomejen sebiadjungiran operator, 
%tako da je njegova resolventa $(D + i)^{-1}$ kompaktna in je $[D, a] = Da - aD$ omejen operator za vsak $a \in \mathfrak{A}$ (DOPOLNI).
%Če s $P$ označimo projekcijo na lastne podprostore $D$, potem želimo strukturo spektralne trojice $(\mathfrak{A}, \mathcal{H}, D)$
%izluščiti s pomočjo trojic $(P \mathfrak{A} P, P \mathcal{H}, PDP)$, toda $P\mathfrak{A}P$ v tem primeru ni več $*$-algebra,
%temveč operatorski sistem 
(za več o tem pristopu glej \cite{connes2021}), najdemo jih pa tudi v kvantni teoriji informacij kot \emph{kvantne grafe} (\cite{weaver2021}, \cite{duanseveriniwinter2013}).
Osnovne lastnosti operatorskih sistemov obravnavamo v poglavju \ref{sec:3}.

Operatorski sistemi so poseben primer operatorskih prostorov,
ki so enotski (ne nujno sebiadjungirani) podprostori $C^*$-algeber.
Vsakemu operatorskemu prostoru $\mathcal{M} \subseteq \mathfrak{A}$ 
pa lahko priredimo operatorski sistem $\mathcal{M} + \mathcal{M}^*$.
Morfizmi na operatorskih prostorih bodo povsem omejene preslikave.
Preplet povsem omejenih preslikav in povsem pozitivnih preslikav bomo obravnavali v poglavju \ref{sec:4}.
Glavna rezultata tega razdelka bosta, da so vse povsem pozitivne preslikave povsem omejene 
in da lahko vsako enotsko povsem skrčitev na operatorskem sistemu $\mathcal{M}$
razširimo do povsem pozitivne preslikave na operatorskem sistemu $\mathcal{M} + \mathcal{M}^*$.
Od 1980-ih let naprej so Ruan, Paulsen, Wittstock in ostali avtorji podrobneje razvijali
teorijo popolnoma omejenih preslikav in operatorskih prostorov, ki jo 
lahko razumemo kot tesen analog teorije popolnoma pozitivnih preslikav in operatorskih sistemov.
V tem magistrskem delu se s tem ne bomo ukvarjali; več podrobnosti lahko bralec najde v \cite{paulsen2002} ali
\cite{pisier2003}.

V krajšem razdelku \ref{sec:5} bomo nato dokazali Stinespringovo karakterizacijo povsem pozitivnih 
preslikav na komutativnih $C^*$-algebrah. Povsem pozitivne preslikave $\varphi: \mathfrak{A} \to \mathfrak{B}$,
kjer je vsaj ena izmed $C^*$-algeber $\mathfrak{A}, \mathfrak{B}$ komutativna, imajo posebno lep opis; 
sovpadajo namreč z navadnimi pozitivnimi preslikavami.

Sledi poglavje \ref{sec:6} o dilatacijah.
Teorija dilatacij, ki sta jo začela Sz.-Nagy in Foias v 1950-ih letih (\cite{nagyfoias2010}),
se okvirno ukvarja s tem, kdaj lahko 
operator $T$ na Hilbertovem prostoru $\mathcal{H}$ realiziramo kot $P_{\mathcal{H}} S |_{\mathcal{H}}$,
kjer je $S$ operator na večjem Hilbertovem prostoru $\mathcal{K} \subseteq \mathcal{H}$
in $P_{\mathcal{H}}$ ortogonalna projekcija s $\mathcal{K}$ na $\mathcal{H}$.
Ponavadi želimo, da ima operator $S$ (ki mu pravimo \emph{dilatacija} $T$)
še nekatere dodatne lastnosti: obravnavali bomo Sz.-Nagyjev zgled,
ki nam pove, da lahko vsak skrčitven operator dilatiramo do izometrije
in vsako izometrijo do unitarnega operatorja. 
Glavni izrek tega poglavja je že zgoraj omenjeni Stinespringov izrek, ki povezuje teorijo dilatacij s povsem pozitivnimi preslikavami.
Dilatacije srečamo med drugim v mehaniki,
kjer lahko evolucijsko polgrupo ireverzibilnega sistema pod določenimi pogoji dilatiramo do unitarne evolucijske grupe (večjega) reverzibilnega sistema 
(glej recimo \cite{MielkePeletierZimmer2025}).

Arvesonovem delu iz leta 1969 sta zgodovinsko sledila Krausov članek \cite{kraus1971} iz leta 1971 
ter Choijev članek \cite{choi1975} iz leta 1975, v katerem sta avtorja klasificirala povsem pozitivne preslikave $\varphi: M_n (\C) \to M_m(\C)$.
Njuno delo med drugim povzamemo v poglavju \ref{sec:7}, v katerem obravnavamo povsem pozitivne preslikave $\varphi: \mathfrak{A} \to \mathfrak{B}$,
kjer je vsaj ena izmed $C^*$-algeber $\mathfrak{A}, \mathfrak{B}$ oblike $M_n (\C)$. 
Posebnega pomena bo trditev, ki pravi, da lahko povsem pozitivno preslikavo $\varphi: \mathcal{S} \to M_n (\C)$
na poljubnem operatorskem sistemu $\mathcal{S} \subseteq \mathfrak{A}$ razširimo do povsem pozitivne preslikave na celotni $C^*$-algebri $\mathfrak{A}$.
V poglavju \ref{sec:8} bomo to trditev dokazali še za povsem pozitivne preslikave $\varphi: \mathcal{S} \to \bh$,
s čimer bomo dokazali Arvesonov razširitveni izrek.

Naslednji pomemben članek za teorijo povsem pozitivnih preslikav sta leta 1977 napisala Choi in Effros (\cite{choieffros1977}),
v katerem sta avtorja dokazala, da imajo vsak operatorski sistem abstraktno karakterizacijo, ki je torej neodvisna od konkretne $C^*$-algebre,
v kateri je operatorski sistem vložen. Choi-Effrosevo abstraktno karakterizacijo obravnavamo v razdelku \ref{sec:9}.
Zato se poraja vprašanje, ali obstaja $C^*$-algebra, v katero lahko vložimo abstrakten operatorski sistem, in ki je na nek način `minimalna'.
Ekvivalentno vprašanje lahko postavimo za konkreten operatorski sistem $\mathcal{S} \subseteq C^*(\mathcal{S})$: ali obstaja minimalna $C^*$-algebra 
$\mathfrak{A}$, tako da obstaja enotska povsem izometrija $i: \mathcal{S} \to \mathfrak{A}$ in je $\mathfrak{A} = C^*(i(\mathcal{S}))$?
Takšni $C^*$-algebri $\mathfrak{A}$ bomo rekli \emph{$C^*$-ogrinjača}.

Problem obstoja $C^*$-ogrinjače je obravnaval že Arveson v \cite{arveson1969}.
V ta namen je vpeljal nekomutativen analog geometrijskega pojma robu Šilova, ki smo ga omenili na začetku uvoda.
Če je $\mathcal{S} \subseteq C(X)$ operatorski sistem, ki loči točke na kompaktnem Hausdorffovem kompaktu $X$,
potem je njegov rob Šilova definiran kot najmanjša zaprta podmnožica $Y \subseteq X$,
tako da vsaka funkcija iz $\mathcal{S}$ na $Y$ doseže svojo supremum normo.
Rob Šilova izhaja iz preučevanja \emph{uniformiranih algeber}:
to so enotske podalgebre (ne nujno zaprte za $*$) $\mathcal{A} \subseteq C(X)$, ki {ločijo točke na $X$}
(za več o tem glej \cite{gamelin1969}). Da lahko vpeljemo nekomutativen rob Šilova,
moramo njegovo definicijo prevesti v besednjak $C^*$-algeber.
To, da operatorski sistem $\mathcal{S} \subseteq C(X)$ loči točke na $X$,
je po Stone-Weierstrassovem izreku ekvivalentno temu, da $\mathcal{S}$ generira $C^*$-algebro $C(X)$.
Ker so zaprte podmnožica $X$ v bijektivni korespondenci z zaprtimi ideali v $C(X)$, lahko vpeljemo 
nekomutativen rob Šilova za poljuben operatorski sistem $\mathcal{S} \subseteq C^*(\mathcal{S})$ kot zaprt ideal $I \subseteq C^*(\mathcal{S})$. 
Arveson je opazil, da če za operatorski sistem $\mathcal{S} \subseteq C^*(\mathcal{S})$ obstaja $C^*$-ogrinjača $C^* _{\env} (\mathcal{S})$,
potem obstaja tudi njegov ideal Šilova $I \subseteq C^*(\mathcal{S})$ in velja $\quot{C^*(\mathcal{S})}{I} \cong C_{\env} ^* (\mathcal{S})$.
V svojem delu je dokazal obstoj ideala Šilova le za določen tip operatorskih sistemov, vendar pa je vprašanje obstoja ideala Šilova in $C^*$-ogrinjače 
za splošen operatorski sistem ostalo nerešeno do
leta 1979, ko je nanj pritrdilno odgovoril Hamana v članku \cite{hamana1979} s svojo teorijo injektivnih ogrinjač.

Arveson je poleg robu Šilova vpeljal še nekomutativen analog tesno povezanega geometrijskega pojma: Choquetovega robu.
Choquetov rob je pojem iz Choquetove teorije, ki se okvirno ukvarja s problemi naslednjega tipa:
\emph{Naj bo $E$ lokalno konveksen prostor in $K \subseteq E$ njegova konveksna, kompaktna podmnožica.
Ali za vsako točko $x \in K$ obstaja verjetnostna mera $\mu$ na $K$, skoncentrirana na $\ext K$, tako da za vsak omejen linearen funkcional $f \in E^*$ velja $\int_X f\, d\mu = f(x)$?}
Pomemben rezultat v tej smeri je Choquetov izrek (\cite[str. 14]{phelps2001}), ki na to vprašanje odgovori pritrdilno, v kolikor je $E$ metrizabilen prostor (več o tej temi v knjigi \cite{phelps2001}).
Ponovno vzamemo operatorski prostor $\mathcal{S} \subseteq C^*(\mathcal{S})$, ki loči točke na $X$. 
Njegov Choquetov rob sestavljajo tiste točke $x \in X$, za katere velja:
če je $\mu$ regularna Borelova verjetnostna mera, tako da je $$f(x) = \int_X f\, d\mu,\quad \forall f \in \mathcal{S},$$
potem je $\mu = \delta_x$. Choquetov rob operatorskega sistema ima to lastnost, da je njegovo zaprtje enako robu Šilova.
Arveson je zato v svojem članku iz leta 1969 vpeljal \emph{robne upodobitve} operatorskega sistema, ki predstavljajo posplošitev točk v Choquetovem robu 
za operatorski sistem funkcij. V tem članku je postavil tudi domnevo, da za poljuben operatorski sistem $\mathcal{S}$ obstaja \emph{dovolj} robnih upodobitev:
to pomeni, da je presek vseh (do unitarne ekvivalence) robnih upodobitev $\mathcal{S}$ enak idealu Šilova.

Kljub temu, da je Hamana dokazal obstoj ideala Šilova, je ta domneva dlje časa ostala nerešena.
Pomemben preboj je nastal leta 2002, ko sta Dritschel in McCollough v članku \cite{dritschel2005}
konstruirala popolnoma nov dokaz obstoja $C^*$-ogrinjače s pomočjo teorije dilatacij. 
Njuno konstrukcijo za operatorske algebre je nato Arveson leta 2003 v neobjavljenem zapisu \cite{arveson2003} posplošil 
na operatorske sisteme. V poglavju \ref{sec:10} predstavimo njihov dokaz, ki mu sledi dokaz obstoja ideala Šilova v poglavju \ref{sec:11}.

Za razliko od Hamanine teorije injektivnih ogrinjač nam argumenti Arvesona, Dritschela in McCollougha omogočajo konstrukcijo robnih upodobitev.
Arveson je nato leta 2008 v članku \cite{arveson2008} dokazal obstoj dovolj robnih upodobitev za separabilne operatorske sisteme. 
Za splošne operatorske sisteme sta trditev dokazala Davidson in Kennedy leta 2015 v članku \cite{davidson2015},
s čimer sta rešila 46-letno domnevo na področju operatorskih algeber.
V zadnjem poglavju \ref{sec:12} predstavimo njuno rešitev.






